\section{智能体评估基准}

\subsection{MiniWoB：简单沙盒环境}

MiniWoB 是最简单的智能体评估环境，提供模拟的网页交互任务（如在模拟 Twitter 上转发、在模拟邮件客户端转发邮件等）。

特点：
\begin{itemize}
    \item 非真实网站，仅为沙盒模拟
    \item 短周期任务——大多数任务只需 $\leq 3$ 个动作
    \item 尽管如此简单，最好的语言模型在 zero-shot 设置下仍\textbf{远未达到完美}
\end{itemize}

\subsection{WebArena：近真实环境}

WebArena 更接近真实世界：

\begin{itemize}
    \item 沙盒环境但高度模拟真实网站（电商、社交媒体、地图等）
    \item 支持\textbf{多标签浏览}——智能体需要在多个应用间切换
    \item 评估\textbf{功能正确性}——判断最终结果是否正确，而非动作序列是否与预设一致
\end{itemize}

\subsection{WebLinx：真实网站交互}

WebLinx 在真实网站上进行评估，增加了一个新的动作：\textbf{与用户沟通}——智能体可以请求缺失信息（如信用卡号）。但它不是一个交互环境，仅是一个交互记录集合，因此无法进行探索或在线学习。

\begin{figure}[H]
\centering
\includegraphics[width=0.95\textwidth]{slides-images/slide-054.jpg}
\caption{WebArena 环境示例：模拟真实网站的沙盒环境，包含电商、社交媒体和工具类应用\protect\footnotemark}
\end{figure}
\footnotetext{Slides 第55页。}

\subsection{本章小结}

从 MiniWoB 到 WebArena 再到 WebLinx，评估基准越来越接近真实世界：任务更复杂、环境更真实、交互更多样。但在所有基准上，语言模型的表现仍远低于人类水平，尤其是在需要多步规划的任务上。

%% ========== Part 9: 合成训练数据 ==========

